\documentclass[11pt,a4paper]{article}

% ======================================================
% 基础宏包
% ======================================================
\usepackage[UTF8,fontset=mac]{ctex}
\setCJKmainfont{Songti SC}
\setCJKsansfont{Heiti SC}
\usepackage{amsmath,amssymb,amsthm,bm}
\usepackage{geometry}
\usepackage{hyperref}
\usepackage{enumitem}
\usepackage{cite}
\usepackage{xcolor}
\usepackage{titlesec}
\usepackage{gbt7714}

% ======================================================
% 页面与样式
% ======================================================
\geometry{margin=1.1in}

\definecolor{linkblue}{RGB}{0,51,153}
\definecolor{citegreen}{RGB}{0,120,80}
\definecolor{urlblack}{RGB}{30,30,30}

\hypersetup{
    colorlinks=true,
    linkcolor=linkblue,
    citecolor=citegreen,
    urlcolor=urlblack,
    pdfauthor={刘晓华},
    pdftitle={分布式认知动力学的物理极限}
}

\titleformat{\section}
  {\Large\bfseries}
  {\thesection}
  {1em}
  {}

\titleformat{\subsection}
  {\large\bfseries}
  {\thesubsection}
  {1em}
  {}

% ======================================================
% 定理环境（定律级）
% ======================================================
\newtheorem{law}{定律}
\newtheorem{theorem}{定理}
\newtheorem{proposition}{命题}
\newtheorem{definition}{定义}
\newtheorem{conjecture}{猜想}

% ======================================================
% 数学符号
% ======================================================
\newcommand{\thetaGT}{\theta^{\star}}
\newcommand{\syn}{\Delta_{\mathrm{syn}}}
\newcommand{\config}{\bm{h}}
\newcommand{\manifold}{\mathcal{M}}
\newcommand{\proj}{\mathcal{P}}
\newcommand{\protocol}{\Pi}

% ======================================================
% 标题信息
% ======================================================
\title{\Huge\bfseries
分布式认知动力学的物理极限\\[0.6em]
\Large
——碎片化观测、隐私边界与协同知识的不可恢复性
}

\author{\Large 刘晓华}
\date{\Large \today}

\begin{document}
\maketitle

% ======================================================
\begin{abstract}
\noindent
在隐私保护与物理隔离成为常态的分布式智能系统中，
协同知识是否\emph{在物理上可被恢复}仍缺乏统一理论。
本文提出“分布式认知动力学”（Distributed Cognitive Dynamics, DCD）
作为一项长期研究纲领，
旨在刻画碎片化观测、隐私边界与能量约束下
协同信息的可达极限。
我们提出三条基本定律与两项不可恢复性定理，
指出线性聚合协议在一类本质协同任务上具有结构性失效，
并引入“信息视界”与“协同容量”概念，
将分布式学习问题提升为
受热力学与信息论双重约束的自然定律问题。
\end{abstract}

% ======================================================
\section{引言：从算法改进到自然定律}
% ======================================================

自 Shannon 建立通信容量理论以来，
信息处理逐渐被视为一类受自然定律约束的过程~\cite{shannon1948}。
在计算物理层面，
Landauer 与 Bennett 进一步揭示了
信息擦除与能量耗散之间不可规避的联系~\cite{landauer1961irreversibility,bennett1982thermodynamics}。

近年来，联邦学习与隐私计算试图在
\emph{不共享原始数据}的前提下实现协同建模~\cite{mcmahan2017communication}。
然而，现有研究大多关注算法收敛性、
通信效率与经验性能，
而较少触及如下根本问题：

\begin{quote}
\emph{在隐私边界与碎片化观测同时存在的条件下，
协同知识是否在物理与信息论意义上仍然可恢复？}
\end{quote}

本文的核心主张是：
在某些条件下，
协同信息并非“难以学习”，
而是\textbf{在自然意义上不可恢复}。
这种不可恢复性并不依赖于模型规模、
优化策略或训练技巧，
而是由信息投影、隐私阻尼与能量耗散共同决定。

% ======================================================
\section{问题设定与协同信息}
% ======================================================

设存在 $N$ 个主体，
每个主体仅能观测真相 $\thetaGT$ 在其观测子空间上的投影：
\[
x_i = \proj_i(\thetaGT), \quad i=1,\dots,N.
\]

\begin{definition}[协同信息]
基于偏信息分解（Partial Information Decomposition, PID）理论~\cite{williams2010nonnegative}，
我们将协同信息 $\syn$ 定义为：
仅在联合观测中出现，
且无法由任何单方或子集线性重构的信息成分。
\end{definition}

典型实例包括 XOR / parity 类任务，
其中任意单方观测与标签条件独立，
但联合观测可完全确定真相。
这类任务构成了研究协同不可恢复性的最小反例族。

% ======================================================
\section{三条基本定律}
% ======================================================

\begin{law}[协同容量定律]
在给定隐私强度 $\epsilon$、
通信预算 $B$ 与能量预算 $W$ 的条件下，
任意分布式协议 $\protocol$
可提取的协同信息满足上界：
\[
\syn(\protocol)
\;\le\;
\mathsf{Cap}_{\mathrm{syn}}(\epsilon,B,W;\{\proj_i\}).
\]
\end{law}

该定律将协同学习问题
从算法设计提升为容量问题，
其地位类似于 Shannon 信道容量~\cite{shannon1948}。

\begin{law}[信息视界定律]
存在信息视界函数 $\mathcal{H}$，
使得当
\[
\gamma(\epsilon) > \mathcal{H}(B,W,N),
\]
系统中协同信息在统计意义下不可识别，
即 $\syn \rightarrow 0$。
\end{law}

该现象体现为
\textbf{认知相变}，
类似物理系统中的相干—退相干转变。

\begin{law}[非线性相干必要性定律]
若协议 $\protocol$ 的全局更新算子
仅依赖于各方消息的一阶线性组合，
则在一类本质协同任务族上，
其协同容量严格为零。
\end{law}

% ======================================================
\section{不可恢复性定理}
% ======================================================

\begin{theorem}[碎片化协同不可恢复定理]
存在任务族 $\mathcal{T}_{\mathrm{syn}}$
与观测投影族 $\{\proj_i\}$，
满足：
\begin{enumerate}[leftmargin=1.6em]
\item 对任意主体 $i$，有 $I(Y;X_i)=0$；
\item 联合观测下 $\syn>0$；
\item 对任意满足隐私约束且线性可分的协议 $\protocol$，
协同信息的估计误差存在严格下界。
\end{enumerate}
\end{theorem}

该定理表明，
线性聚合在协同任务上的失败是
\textbf{结构性且不可避免的}，
与模型复杂度或优化技巧无关。

\begin{theorem}[隐私—能量协同上界定理]
在 Markov 毯约束~\cite{friston2010free}
与 Landauer 原理~\cite{landauer1961irreversibility,bennett1982thermodynamics}下，
任意协议对协同信息的增益满足：
\[
\syn
\;\le\;
\min\!\left\{
I(\thetaGT;B_{1:N}),
\;\frac{W_{\mathrm{diss}}}{kT\ln 2}
\right\}.
\]
\end{theorem}

该结果表明，
分布式学习能力直接受热力学第二定律制约。

% ======================================================
\section{可达机制：最小非线性相干}
% ======================================================

不可恢复性并不否认协同的可能性，
而是限定了其\emph{必要结构条件}。

恢复协同所需的最小机制是
对交互统计的显式建模，
例如二阶相干矩：
\[
C_{ij}
=
\mathbb{E}[z_i \otimes z_j]
-
\mathbb{E}[z_i]\otimes\mathbb{E}[z_j].
\]

该思想与信息场论中
对高阶相关的系统建模一致~\cite{ensslin2009information}，
并构成张量相干提取
（Tensor Coherence Extraction, TCE）的理论基础。

% ======================================================
\section{研究路线与长期愿景}
% ======================================================

本文并非提出单一算法，
而是试图奠定一套关于
\textbf{分布式知识可达性的自然定律体系}。

短期研究将聚焦于：
\begin{itemize}
\item 不可恢复性定理的严格化证明；
\item 协同容量上界的可计算表达；
\item 多中心医疗与社会协作系统中的信息视界实证。
\end{itemize}

从长远看，
该理论有望为群体智能、生物协同与社会认知
提供统一的物理与信息论基础。

\bibliographystyle{gbt7714-numerical}
\bibliography{refs}

\end{document}
